% !TeX root = main.tex
\section{定位与基于位置的服务}

2005年到2014年间，德国的LBS服务提供商数量经历了爆发式增长，年增长率一度非常高。不过这一波增长主要是消费市场驱动的，比如地图导航、社交打卡、外卖跑腿这些面向普通消费者的应用。相比之下，工业生产场景里对LBS的应用还比较有限，目前主要集中在物流环节，比如仓库里货物和叉车的追踪。消费领域已经把这项技术玩得很成熟了，但工业领域还有很大的应用空间没有被打开。

\subsection{原理与方法}

常见的定位原理分为以下几大类。
\begin{enumerate}
    \item \textbf{信号强度法（Signalstärke）}：信号在传播过程中会随距离增大而衰减，只要提前测出一张"信号强度地图"，就能反推出接收点大致在哪里，GSM、UMTS、WLAN、蓝牙定位都属于这个思路，缺点是容易受到障碍物遮挡和信号反射的干扰。
    \item \textbf{飞行时间测量（Laufzeitmessung）}，也就是通过信号传播的时间差来算距离，光速是已知的（大约100皮秒对应3厘米），所以只要测出时间就能算出距离。这里面又细分为\textbf{到达时间法（TOA, Time of Arrival）}、\textbf{到达时间差法（TDOA, Time Difference of Arrival）}和\textbf{往返时间法（RTT, Round Trip Time）}，GPS、伽利略、格洛纳斯这些卫星导航系统用的都是这一类原理。
    \item \textbf{三边测量与三角测量（Triangulation）}，前者是通过测量到多个参照点的距离来确定位置，后者是通过测量角度，立体视觉、光切法、WLAN和蓝牙的一些定位方案会用到这类思路。
    \item \textbf{特征匹配法（Merkmalsbasiert）}：把传感器测到的环境特征（比如激光雷达扫出来的点云）和已有的地图做匹配，比较典型的做法是ICP迭代最近点算法，不过这种方法需要一个比较好的初始猜测值才能收敛。
    \item \textbf{惯性传感法（Intertialsensorik）}：靠加速度计、陀螺仪、地磁传感器等器件，从一个已知的起点开始不断累加位移，缺点是误差会随时间累积，所以实际应用中经常和其他方法组合使用，也就是常说的"混合定位"。
    \item \textbf{相位测量法（Phasenmessung）}：通过测量信号在一个周期内的相位来确定相对或绝对位置，比如利用干扰信号计数脉冲，这种方法精度较高，但通常需要配合其他手段一起使用。
\end{enumerate}
在实际选型时，需要沿着若干维度对候选系统做综合权衡：\textbf{精度（Genauigkeit）}是否满足任务要求；\textbf{作用范围（Arbeitsraum）}能否覆盖目标工作空间；\textbf{基础设施需求（Infrastrukturbedarf）}有多大，是否必须预先布设锚点、反射板或轨道；\textbf{抗干扰能力与鲁棒性（Robustheit）}在存在遮挡、反射和金属结构的车间环境中是否足够；\textbf{可扩展性（Skalierbarkeit）}以及可同时定位的\textbf{目标数量（Teilnehmeranzahl）}能否随产线规模增长；\textbf{可用性（Verfügbarkeit）}，即系统在多大比例的时间和空间内能给出有效位置；所用\textbf{频段（Frequenzen）}是否合规、是否与厂内既有无线系统冲突；\textbf{成本（Kosten）}，包括设备、安装与运维；\textbf{信息安全（Security）}；以及该技术的\textbf{市场成熟度（Marktbedeutung）}与供应商生态。图  直观地展示了各种室内定位技术在精度和作用范围这两个维度上的分布，从毫米级的摄像头、超宽带（UWB）定位，到覆盖整栋建筑的WLAN和蜂窝网络定位。

\begin{figure}[htbp]
  \centering
  \includegraphics[width=0.7\textwidth]{indoorlokalisierungverteilung.png}
  \caption{各种室内定位技术在精度和作用范围这两个维度上的分布}
  \label{fig:indoorlokalisierungverteilung}
\end{figure}

目前室内定位领域还没有出现像GPS在室外那样一统天下的主流方案，这也是这个领域仍然活跃、值得研究的原因之一。

\subsection{地图}

有了定位能力之后，接下来要解决的问题是这些位置信息怎么组织和呈现，也就是地图。例如我现在要找到最近的服装店或餐厅，并从我当前位置导航过去，这个简单需求背后其实包含了四层能力：把信息画出来给人看（图形化展示）、按条件搜索兴趣点（POI搜索）、规划一条路线、以及实时导航。

地图本质上是数据，只是被组织成了人类容易理解的形式。不同的应用场景需要不同的数据结构来支撑，课程介绍了室内环境建模的四种典型思路。
\begin{enumerate}
    \item \textbf{基于网格的几何建模（Cell based）}，也叫占据栅格地图，把空间切成一个个小格子，每个格子标记是否被占用。
    \item \textbf{基于边界的几何建模（Boundary Based）}，用轮廓线来描述空间和物体的边界。
    \item \textbf{基于符号或地点的建模（Symbolic/Place based）}，不追求精确的几何形状，而是用"房间""通道"这类语义化的地点概念来组织空间。
    \item \textbf{混合式建模（Hybrid）}，把面向对象的方法和边界建模结合起来，兼顾精确性和语义信息。
\end{enumerate}
这几种方式各有取舍，选哪一种取决于具体应用对精度、灵活性和计算成本的要求。



\subsection{基于位置的服务}

有了定位能力和地图这两块基础之后，就可以在其上构建各类基于位置的服务（LBS）。汇总出以下十种类型：
\begin{enumerate}
  \item \textbf{参照定位（Referenzierung）}：为后续一切位置信息提供一个基准，即把对象的位置表达为相对于某个已知参考系或参考点的相对位置，本身不关心姿态。
  \item \textbf{定位（Positionierung）}：在参照的基础上确定对象的实际位置，通常不只给出位置，还包括姿态，即完整的位姿（Pose）；这是与参照定位的主要区别。
  \item \textbf{导航（Navigation）}：在已知自身位姿和环境模型的前提下，给出应当行驶的路径或轨迹，并在行驶过程中持续跟踪与修正，是移动机器人和搬运车辆的核心服务。
  \item \textbf{质量保证（Qualitätssicherung）}：借助位置信息保证工序在正确的位置、按正确的顺序执行，重点在于避免错误的产生，而不是事后把错误检出来，例如拧紧工具只有在正确的螺栓位置上才被使能。
  \item \textbf{文档记录（Dokumentation）}：对位置信息做持续的追踪和留痕，往往是出于法规或合规要求（如可追溯性、责任认定），记录本身即是目的。
  \item \textbf{数据分析（Analyse）}：在文档记录所积累的位置数据之上做进一步的评估，比如识别搬运路径中的绕行与瓶颈、统计设备利用率，用于流程优化；与文档记录的区别在于目的是分析而非留证。
  \item \textbf{信息提供（Information）}：依据使用者当前所处的位置提供相应的信息，比如给工人推送与他当前工位、当前任务相关的作业指导，这也是增强现实（AR）应用的基础。
  \item \textbf{安全防护（Safety）}：防止人员受到设备的伤害，比如人员接近危险区域时自动预警、降速或停机。
  \item \textbf{防范恶意行为（Security）}：也包括厂区与资产的防护，识别人或设备的异常行为并采取应对措施；由于两者在技术手段上高度重叠，Security 与 Safety 之间的界限在实践中并不容易划清。
  \item \textbf{绩效管控（Leistungssteuerung）}：借助位置数据确定所需的绩效参数（如节拍、行程、等待时间），并在必要时对人员或设备的工作表现加以调控。
\end{enumerate}

值得一提的是，定位技术在工业场景中的应用是具有许多风险和限制的，尤其是当定位信息被用来监控员工的行为和绩效时。比如键盘输入速度监测、通话记录、屏幕操作追踪，这些技术手段虽然可行，但会直接触及员工的人格权和隐私权，在德国的实际操作中，这类系统的引入和使用受到《联邦数据保护法》的约束，并且必须经过企业职工委员会的共同决策程序，不能由企业单方面决定。

真实案例：运动软件Strava发布的全球运动热力图，因为记录了大量军人在军事基地内跑步的轨迹，无意中暴露了一些本应保密的军事设施位置和内部路网布局，这个例子很直观地说明了定位数据一旦被聚合和分析，可能带来意想不到的信息泄露风险。这提醒我们，定位技术在带来便利的同时，也需要认真考虑数据安全和个人隐私的边界。

\paragraph{示例：厂内物流中的定位应用}

在如图  所示的这样一个工厂中，有许多自动搬运机器人，它们要完成货物的取放和运输任务，需要具备几项核心能力：
\begin{enumerate}
  \item \textbf{任务规划（Handlungsplanung）}：把一个大任务拆解成一步步可执行的动作
  \item \textbf{导航（Navigation）}：包括自身定位、路径规划和路径跟踪
  \item \textbf{避障（Kollisionsvermeidung）}：绕开障碍物、识别危险情况
  \item \textbf{物体识别（Objekterkennung）}：认出托盘、行人、其他车辆
  \item \textbf{交通调度（Verkehrsregelung）}：多台机器人之间的优先级协调和死锁避免
\end{enumerate}

\begin{figure}[htbp]
  \centering
  \includegraphics[width=0.7\textwidth]{intralogistik.png}
  \caption{需要许多自动搬运机器人的工厂}
  \label{fig:intralogistik}
\end{figure}

导航环节里的自身定位，对比了几种在物流仓库里常见的定位方案：光学或电磁轨道引导（车辆沿着地面铺设的轨迹行驶）、磁点导航（依靠预埋在地面的磁性标记点）、二维激光反射板定位（依靠悬挂在环境中的反射板作为参照）、以及基于自然地标的定位（不需要额外布置标记物，直接利用环境本身的特征，又分为二维和三维两种做法）。这几种方案在地图生成的难易程度、定位精度、对生产流程调整的适应能力、以及安装部署的成本上各有优劣。

\paragraph{示例：航空维修行业的实际项目}

第二个案例来自航空维修领域，讲的是汉莎技术公司（Lufthansa Technik）的"GSE 4.0"项目。GSE指的是机场地面保障设备（Ground Support Equipment），如图 \ref{fig:gse40} 所示，汉堡基地目前有超过1500台这样的设备在使用，而一架飞机在地面停留维修期间，大约需要用到500台GSE。

\begin{figure}[htbp]
  \centering
  \includegraphics[width=0.9\textwidth]{gse40.png}
  \caption{Lufthansa Technik 的 GSE 4.0 项目}
  \label{fig:gse40}
\end{figure}

这个项目要解决的问题很实际：设备使用完之后经常被随手放在没有固定规划的地方，下一个需要用它的机修工人就得花时间到处寻找，白白浪费了工时，也没法把设备的利用率提到最高。项目的解决思路是给每台设备装上定位标签，让设备之间通过蓝牙（BLE）互相定位，具体的定位方式是基于信号强度（RSS）；同时在建筑物上安装若干锚点标签作为本地定位的参照物；设备的标签在设定的事件触发时把自己的位置信息发送给LoRa网关（这种网关的信号覆盖范围可以达到几公里）；网关再通过移动网络把位置数据统一传回中心数据库。有了这套基础设施之后，就可以在此之上搭建出好几种实用的LBS功能：设备的实时追踪和状态查询、按需的临时预约申请、以及提前的使用计划安排。最终效果是让设备管理从原来分散、被动的状态，变成\textbf{集中协调、可以提前规划的状态，明显提高了透明度，也减少了设备闲置或找不到的时间}。

\paragraph{工业室内定位}

IFPT研究所正在做的科研项目 IIL（Industrielle Indoor Lokalisierung）。

矩阵图 \ref{fig:iil} 说明了这个项目的出发点：生产制造过程中的不同环节，比如寻找物料、搬运、抓取、组装、加工、检测——对定位精度的要求差异非常大，从米级（找东西大概知道在哪个区域就够了）一直到微米级（比如精密测量环节），这些不同层级的需求分别对应着信号强度、飞行时间、三边测量、三角测量等不同的技术手段，没有一种定位技术能够单独覆盖所有场景。

\begin{figure}[htbp]
  \centering
  \includegraphics[width=0.9\textwidth]{iil.png}
  \caption{生产制造过程中不同层级的需求}
  \label{fig:iil}
\end{figure}

用大家都很熟悉的谷歌地图做类比，一个好的环境模型应该是分层和可集成的：卫星图像有不同精细程度的层级，再叠加街景图、社区标注、路况信息、矢量数据等，这些不同来源、不同类型的数据整合在一起，才构成了我们平时用起来很顺手的地图产品。IIL项目想做的，就是把工厂里各种来源的定位信息（比如移动机器人的定位、增强现实设备的定位、智能工具的定位、移动测量设备的定位）也用类似的思路整合起来，形成一个统一的位置信息底座，再向上支撑数字工厂、LBS等各类应用。